\chapter{Cable modelization}

\section{Introduction}\label{cable_modelization_sec_introduction}

Remotely Operated Vehicles (ROVs) are used to explore unknown areas. Their tether is the only way
to provide direct feedback to the pilot, as electromagnetic waves are absorbed in the first few
centimetres of water. However, the tether can interfere with the robot's movement and, if its shape
is not controlled, it can become entangled in obstacles. This is even more critical when exploring
confined karstic terrain, wrecks or offshore installations with small observation ROVs. This study
focuses on shape estimation of tethers of a few metres in length, \textit{e.g.} the portion of
cable connecting two mini-ROVs of a chain~\cite{laranjeira2020a} or a mini-ROV to an Unmanned
Surface Vehicle (USV)~\cite{tortorici2019}. To be used in robot control for path following and
obstacle avoidance, the chosen geometric model must be parametrised and fast to compute.

The standard catenary model, bound to a vertical plane, is defined as the shape of an idealized
homogeneous hanging cable with fixed length and fixed ends, only subject to its own weight in the
air. The catenary shape can be determined knowing the cable’s length and its deflection. Depending
on the available sensors, its shape can be estimated either from the relative position of the two
attachment points such as in~\cite{dantonio2021, DAntonio2022} or from the local cable tangent
orientation~\cite{Dru_Dun2022b}.

However, the extension of the catenary model to an underwater cable is not straightforward in the
presence of hydrodynamic effects. In static conditions, submarine cables are only subject to weight
and buoyancy. While the quasi-static assumption remains applicable for aerial cables at low speed
when air friction can be neglected, the underwater drag involves hydrodynamic forces that move the
cable from its vertical plane, even at low speed~\cite{Dru_Dun2022a}.

The main contribution of the present study is the introduction of a new enhanced catenary model
that incorporates full tilting of the tether during surge or sway motion of the underwater robot.
The relevance of this geometric model is confirmed by experiments conducted on eight cables with
different mechanical properties (\cref{cable_modelization_fig_experimentalsetup}).

\begin{figure}
	\centering
	\includegraphics[width=\figurewidth, trim=40 20 110 20, clip]{experiment_picture.jpg}
	\caption{Experimental setup. A dark red cable connects the ROV to a fixed point. Optical markers are regularly placed on
		the cable for underwater motion tracking. The slack light yellow tether is used for communications with the ROV}
	\label{cable_modelization_fig_experimentalsetup}
\end{figure}

The paper is organized as follows. \cref{cable_modelization_sec_related_work} presents related work
on underwater tethered robots and cable shape estimation.
\cref{cable_modelization_sec_modelization} introduces the augmented catenary model. The
experimental protocol and data analysis are detailed in \cref{cable_modelization_sec_experiment}.
\Cref{cable_modelization_sec_conclusion} concludes the study and outlines future work.

\section{Related work}
\label{cable_modelization_sec_related_work}

Tethered robots are particularly common in underwater robotics, where cables are the only way to
provide real time communication between ROVs and the surface~\cite{tortorici2019}. Multiple
underwater robots can also be connected together in order to communicate freely and share
resources. In~\cite{yu2013}, small vehicles are connected to a bigger autonomous submarine. In
addition, underwater robot chains are being investigated in order to limit the mechanical impact of
tethers on ROV navigation and control by the addition of multiple mini-ROV on the cable to control
its shape~\cite{laranjeira2020a, Dru_Dun2022b}. It results in an additional underwater tethered
robot configuration in which small devices are linked together by short cable portions.

The safe control of tethered robotic systems strongly relies on cable shape knowledge: on the one
hand, controlling tethered robots requires knowledge of the cable's state~\cite{laranjeira2020a},
on the other hand, cable shape knowledge gives information about the localization of the robot at
its end point~\cite{mcgarey2017, Viel2023}. Some underwater cables are designed as proprioceptive
sensors in order to measure their own shape. Optical-fiber cables measure their deformations using
reflectometry techniques~\cite{duncan2007, floris2021}. Such cables can be used to locate connected
underwater vessels~\cite{yu2013}. Another version of optical fiber cable shape sensing consists in
adding an external optical fiber sensor coating on a cable~\cite{xu2016}. Such cables are, however,
fragile and extremely expensive. The cable shape can also be estimated from the measurements of
inertial sensors placed along it~\cite{franck2013}. However, these methods require the use of a
specially designed cable. In addition, the accuracy of these methods decreases with the cable's
length due to error accumulation, e.g.\ approx.\ \qty{1}{\cm} error for a \qty{1}{\m} long optical
fiber cable~\cite{duncan2007}.

Other strategies make use of physical or geometrical cable models. A physical model can be
developed based on the cable's hydrodynamics and the forces acting upon
it~\cite{eidsvik2018,meng2020,soylu2010,hong2020}. Although such models are the most comprehensive
and designed to closely match the physics, they are computationally intensive and require an
extensive understanding of multiple parameters which are difficult to measure under actual
circumstances, such as water current or thruster parameters. As a result, simpler models are often
preferred. The cable may be constrained to a simplified geometric shape artificially, for instance
by introducing weights or sliding floaters to make it piecewise linear~\cite{Viel2022_oe, Viel2022,
	Viel2023}. Some studies employ the catenary model, a hyperbolic curves, for underwater or aerial
settings~\cite{laranjeira2020a, Dru_Dun2022b}, whereas parabolic curves may also be utilized in
aerial scenarios~\cite{smolentsev_shape_2023}.

However, in the presence of external forces, including ambient currents, hydrodynamic damping
during motion, the geometric models that represent quasi-static cable shapes may not be accurate
enough. A preliminary study~\cite{Dru_Dun2022a} investigated the accuracy of an inclined catenary
model for short cables connecting a pair of moving robots, using motion tracking of genuine cables.
This paper expands on the findings of the study by presenting an enhanced parameterized model that
considers complete cable inclination caused by external forces underwater.

\section{Modeling and estimation}
\label{cable_modelization_sec_modelization}

Let $\threedframe{w} =(\mathbf{O}_{w}, \mathbf{x}_{w}, \mathbf{y}_{w}, \mathbf{z}_{w})$ be a fixed
Cartesian reference frame with the $\mathbf{z}_{w}$-axis vertical and pointing upwards.

\subsection{Standard catenary model}
\label{cable_modelization_subsec_catenary_model}

The shape of a homogeneous hanging cable, with fixed length $L$, only subjected to its weight and
buoyancy, is defined by a standard catenary curve
(\cref{cable_modelization_fig_schematic_catenary_model}). The catenary lies in a vertical plane
$\plane{v}=(\mathbf{O}_{v}, \mathbf{x}_{v},\mathbf{z}_{v})$, associated to the Cartesian frame
$\threedframe{v}=(\mathbf{O}_{v}, \mathbf{x}_{v},\mathbf{y}_{v},\mathbf{z}_{v})$ where
$\mathbf{O}_{v}$ is the lowest point of the curve, $\mathbf{x}_{v}$ is horizontal and
$\mathbf{z}_{v}$ is collinear to $\mathbf{z}_{w}$. Curve points coordinates $\left( {}^{v}X,{}^{v}Z
	\right) \in \plane{v}$ are given by:
\begin{gather}
	({}^{v}X, {}^{v}Z)\in[{}^{v}X_{i}, {}^{v}X_{f}]\times \mathbb{R} ~ \text{such that} \nonumber\\
	{}^{v}Z = \frac{\cosh{({}^{v}XC)} - 1}{C}
	\label{cable_modelization_eq_catenary_base}
\end{gather}
where $C \in \mathbb{R}^*_+$ is the catenary parameter, and indexes $i$ and $f$ refer to the curve's endpoints.
The left superscript indicates the frame in which a coordinate is given.

\begin{figure}
	\centering
	\inputtikz{figures/catenary}
	\caption{
		Catenary curve of length $L$ defined in $\plane{v}$.
		Parameters are the sag $H$, the difference of elevation $\Delta H$ and the horizontal distance $l$ between
		attachment points
	}
	\label{cable_modelization_fig_schematic_catenary_model}
\end{figure}

\subsection{Augmented catenary model}
\label{cable_modelization_subsec_augmented_catenary_model}

As soon as the underwater cable is moved fast enough, hydrodynamics forces become non-negligible.
The drag force dampens the movement with an amplitude that is proportional to the square of the
speed. This effect is amplified by the added mass when significant accelerations are
present~\cite{fossen2011, eidsvik2018}.

In the steady state case where all points on the cable move at the same velocity, they are
subjected to the same drag forces. As a consequence, the resulting uniform acceleration of the
cable is tilted, and its intensity is increased depending on the velocity. The addition of two
degrees of freedom in the cable shape model to account for this effect is used here to extend the
standard catenary model.

The first degree of freedom is a rotation of angle $\gamma$ around the catenary axis $\mathbf{e}$,
which is defined as the unit vector that connects cable endpoints~\cite{Dru_Dun2022b}
(\cref{cable_modelization_fig_schematic_gamma_augmented_catenary}). This is to consider the
deformations of a tether subjected to sway motions, \emph{i.e.} the cable ends move out of the
vertical plane, namely $\plane{v}$. This rotation applied to $\threedframe{v}$ and $\plane{v}$
gives $\threedframe{\gamma} = (\mathbf{O}_{\gamma}, \mathbf{x}_{\gamma},\mathbf{y}_{\gamma},
	\mathbf{z}_{\gamma})$ and $\plane{\gamma}$, respectively.

\begin{figure}
	\centering
	\inputtikz{figures/planes}
	\caption{
		Standard (blue) and \gammamodel{} (red) catenary for $\gamma = \qty{10}{\degree}$
	}
	\label{cable_modelization_fig_schematic_gamma_augmented_catenary}
\end{figure}

The second degree of freedom, a rotation of angle $\theta$ in $\plane{\gamma}$
(\cref{cable_modelization_fig_theta_augmented_catenary}), considers tether deformations caused by
surge motions, \textit{i.e.} the two cable ends get closer or further away. The approximation of
uniform hydrodynamic forces along the cable results in a uniform tilted acceleration
$\mathbf{a}_{\theta\gamma}$ with respect to gravity. Let $\threedframe{\theta\gamma} =
	(O_{\theta\gamma}, \mathbf{x}_{\theta\gamma}, \mathbf{y}_{\theta\gamma},
	\mathbf{z}_{\theta\gamma})$ be the associated Cartesian frame where $\mathbf{z}_{\theta\gamma}$ is
parallel to $\mathbf{a}_{\theta\gamma}$ and $\mathbf{y}_{\theta\gamma}$ is parallel to
$\mathbf{y}_{\gamma}$. The two ends of the cable can be expressed in $\threedframe{\theta\gamma}$
and~\cref{cable_modelization_eq_catenary_base} is used to construct the oriented catenary. This
2-DOF model is named the \emph{\thetagammamodel{} model}.

\begin{figure}
	\centering
	\inputtikz{figures/tilted_catenary}
	\caption{
		\gammamodel{} (red) and \thetagammamodel{} (green) catenary for $\theta = \qty{45}{\degree}$
	} \label{cable_modelization_fig_theta_augmented_catenary}
\end{figure}

Note that zero values for $\gamma$ and $\theta$ yield the \emph{standard catenary} model. If
$\theta$ only is set to zero, the model presented in~\cite{Dru_Dun2022b} can be obtained, and is
referenced as \emph{\gammamodel{}}.

\subsection{Curvilinear discretization and residual}

Let us consider $n+1$ 3D measurement points distributed along the cable's length. Accordingly, the
model is discretized with respect to the curvilinear abscissa.

Henceforth, the symbols $m,v, \gamma$ and $\theta\gamma$ are used as indexes referring to the
measurements and the estimates of the \textit{standard}, \gammamodel{} and \thetagammamodel{}
catenary model variations, respectively. Model points are written:
\begin{align*}
	\pestimated{w}{\ast}{k}=({}^{w}X_{k}^{\ast}, {}^{w}Y_{k}^{\ast}, {}^{w}Z_{k}^{\ast}) \\
	k\inzeroton, \ast\in \left\{ m, v, \gamma, \theta\gamma \right\}
\end{align*}
where indexes $k=0$ and $k=n$ are the tether endpoints, also indexed $i$ and $f$.

Model variants accuracies are computed by means of the residuals:
\begin{align}
	\pdistance{\ast} & = \frac{1}{n}\sum_{k=0}^{n}{\norm{\pestimated{w}{m}{k} -\pestimated{w}{\ast}{k}}}, ~
	\ast\in \left\{v, \gamma, \theta\gamma \right\}
	\label{cable_modelization_eq_error_vt_definition}
\end{align}

\subsection{Model parameters estimation}\label{cable_modelization_subsec_estimation}
\subsubsection{Catenary parameter}

Let $L$ be the known tether length, $l$ the horizontal distance and $\Delta H$ the difference in
elevation between the tether attachment points calculated from the measurements of the initial and
final attachment points $\pestimated{\Diamond}{m}{i}$ and $\pestimated{\Diamond}{m}{f}$, with
$\Diamond\in\{v, \gamma, \theta\gamma\}$ the frame in which the points are expressed
(\cref{cable_modelization_fig_schematic_catenary_model}). The catenary parameter $C$ is estimated
by finding the root of function $f$ which relates the curve's length and~$C$~\cite{laranjeira2020a}
(using Brent's method~\cite{brent2013algorithms}):
\begin{align*}
	f(C)=C^2 \left(L^2 - \Delta H^2\right) - 4  \sinh^2(lC/2)
\end{align*}

\subsubsection{Augmented model parameters}

Let $\pdistance{\theta\gamma}\left(\theta, \gamma\right)$ be the accuracy written as a function of
$\theta$ and $\gamma$~\cref{cable_modelization_eq_error_vt_definition}. Angles $\theta$ and
$\gamma$ are estimated by minimizing this function by means of the trust region reflective
algorithm~\cite{Branch1999TIR}:
\begin{align*}
	(\theta, \gamma) = \argmin_{(\theta, \gamma)\in\left[-\pi,\pi\right]^2} \pdistance{\theta\gamma}\left(\theta, \gamma\right)
\end{align*}

The initial estimates $\theta_0$ and $\gamma_0$ are respectively set to zero and to the angle
between $\plane{\gamma_0}$ and $\plane{v}$, where $\plane{\gamma_0}$ is the closest plane to all
measured points.

\section{Experiments}\label{cable_modelization_sec_experiment}

\subsection{Cables description}
\label{cable_modelization_subsec_cable}

The experiments involve eight \qty{3}{\m} long cables presented in
\cref{cable_modelization_fig_image_all_cables}. These cables are meant to reflect the full range of
mechanical characteristics found in robot umbilicals and tethers: stiffness from the most flexible
to the most rigid; neutral, positive or negative buoyancy; elastic or nonelastic
(\cref{cable_modelization_tab_cables_characteristics}). Overall, these cables can be classified
into two categories based on their capacity to retain plastic deformation: cables
\subref{cable_modelization_fig_image_coaxial}, \subref{cable_modelization_fig_image_four_pairs} and
\subref{cable_modelization_fig_image_two_pairs} do, and cables
\subref{cable_modelization_fig_image_floating}, \subref{cable_modelization_fig_image_rope},
\subref{cable_modelization_fig_image_weighted}, \subref{cable_modelization_fig_image_chain} and
\subref{cable_modelization_fig_image_elastic} do not. Cables
\subref{cable_modelization_fig_image_floating}, \subref{cable_modelization_fig_image_weighted} and
\subref{cable_modelization_fig_image_chain}, are expected to give the best results in catenary
shape fitting due to their clear positive or negative buoyancy, highlighted in
\cref{cable_modelization_tab_cables_characteristics}.

\begin{figure}
	\renewcommand\thesubfigure{\arabic{subfigure}}
	\centering
	\begin{subfigure}{.2 \textwidth}
		\centering
		\includegraphics[width=\textwidth]{coaxial.jpg}
		\caption{}
		\label{cable_modelization_fig_image_coaxial}
	\end{subfigure}
	\begin{subfigure}{.2 \textwidth}
		\centering
		\includegraphics[width=\textwidth]{four_pairs.jpg}
		\caption{}
		\label{cable_modelization_fig_image_four_pairs}
	\end{subfigure}
	\begin{subfigure}{.2 \textwidth}
		\centering
		\includegraphics[width=\textwidth]{two_pairs.jpg}
		\caption{}
		\label{cable_modelization_fig_image_two_pairs}
	\end{subfigure}
	\begin{subfigure}{.2 \textwidth}
		\centering
		\includegraphics[width=\textwidth]{floating.jpg}
		\caption{}
		\label{cable_modelization_fig_image_floating}
	\end{subfigure}
	\begin{subfigure}{.2 \textwidth}
		\centering
		\includegraphics[width=\textwidth]{rope.jpg}
		\caption{}
		\label{cable_modelization_fig_image_rope}
	\end{subfigure}
	\begin{subfigure}{.2 \textwidth}
		\centering
		\includegraphics[width=\textwidth]{catenary.jpg}
		\caption{}
		\label{cable_modelization_fig_image_weighted}
	\end{subfigure}
	\begin{subfigure}{.2 \textwidth}
		\centering
		\includegraphics[width=\textwidth]{chain.jpg}
		\caption{}
		\label{cable_modelization_fig_image_chain}
	\end{subfigure}
	\begin{subfigure}{.2 \textwidth}
		\centering
		\includegraphics[width=\textwidth]{elastic.jpg}
		\caption{}
		\label{cable_modelization_fig_image_elastic}
	\end{subfigure}
	\caption{
		Cables ends:
		(\subref{cable_modelization_fig_image_coaxial}) coaxial cable,
		(\subref{cable_modelization_fig_image_four_pairs}) four-pair network cable,
		(\subref{cable_modelization_fig_image_two_pairs}) two-pair network cable,
		(\subref{cable_modelization_fig_image_floating}) floating braided rope,
		(\subref{cable_modelization_fig_image_rope}) sail rope,
		(\subref{cable_modelization_fig_image_weighted}) leaded seamstress rope,
		(\subref{cable_modelization_fig_image_chain}) plain steel chain,
		(\subref{cable_modelization_fig_image_elastic}) elastic
	}
	\label{cable_modelization_fig_image_all_cables}
	\renewcommand\thesubfigure{\alpha{subfigure}}
\end{figure}

\begin{table}
	\centering
	\robustify\bfseries
	\sisetup{ % pour changer le nombre de chiffres après la virgule affichés après compilation
		round-mode = places,
		round-precision = 3}
	\caption{Cables characteristics}
	\begin{tabular}{@{}l|SSSS@{}}
		\toprule
		cable \#                    & \subref{cable_modelization_fig_image_coaxial} & \subref{cable_modelization_fig_image_four_pairs} & \subref{cable_modelization_fig_image_two_pairs} & \subref{cable_modelization_fig_image_floating} \\ \midrule
		diameter / \unit{\m}        & 0.0068                                        & 0.005                                            & 0.0042                                          & 0.0077                                         \\
		weight / \unit{\N}          & 1.0791                                        & 0.6867                                           & 0.3924                                          & 0.2943                                         \\
		weight in water / \unit{\N} & 0.04905                                       & 0.14715                                          & -0.04905                                        & -0.4905                                        \\
		wet weight / \unit{\N}      & 1.0791                                        & 0.6867                                           & 0.3924                                          & 0.981                                          \\
		plastic deformation         & {\checkmark}                                  & {\checkmark}                                     & {\checkmark}                                    & {$\cross$}                                     \\ \midrule[1pt]
		cable \#                    & \subref{cable_modelization_fig_image_rope}    & \subref{cable_modelization_fig_image_weighted}   & \subref{cable_modelization_fig_image_chain}     & \subref{cable_modelization_fig_image_elastic}  \\ \midrule
		diameter / \unit{\m}        & 0.0055                                        & 0.0035                                           & {---}                                           & 0.0054                                         \\
		weight / \unit{\N}          & 1.0791                                        & 1.3734                                           & 3.11958                                         & 0.5886                                         \\
		weight in water / \unit{\N} & 0.24525                                       & 1.1772                                           & 2.70756                                         & 0.0981                                         \\
		wet weight / \unit{\N}      & 1.5696                                        & 1.52055                                          & 3.11958                                         & 0.6867                                         \\
		plastic deformation         & {$\cross$}                                    & {$\cross$}                                       & {$\cross$}                                      & {$\cross$}                                     \\ \bottomrule
	\end{tabular}
	\label{cable_modelization_tab_cables_characteristics}
\end{table}

\subsection{Experimental setup}

The experiment is set up in a \qty{7.2}{\m} long, \qty{4.2}{\m} wide and \qty{3}{\m} deep tank. One
end of the cable, namely $\pestimated{w}{m}{n}$, is attached to a ROV (BlueROV), while the other,
namely $\pestimated{w}{m}{0}$, is fixed and attached to the side of the pool
(\cref{cable_modelization_fig_experimentalsetup}).

The cable's and robot's positions are recorded by a 5-camera, \qty{100}{\Hz}
\emph{Qualisys}\footnote{See specifications at www.qualisys.com/cameras/underwater} underwater
motion capture system. The robot as well as each cable are equipped with passive reflective markers
and the cable's markers are evenly spaced out by \qty{0.2}{\m} ($n=15$). The calibration of the
motion capture system gives a standard deviation of \qty{2.2}{\mm} across the pool's volume for a
known object's length. The points $\pestimated{w}{m}{k}$ are measured in the \textit{Qualysis}
world Cartesian frame which is defined by the positioning of a calibration target that may not be
perfectly vertical in the real experiment. This introduces a small offset
($\pm$\qty{0.05}{\radian}) for the steady state angles as it can be seen in
\cref{cable_modelization_fig_sequences_cable3surge,cable_modelization_fig_sequences_cable3sway,cable_modelization_fig_sequences_cable6surge,cable_modelization_fig_sequences_cable6sway}.

The initial position of the robot is defined such that the robot is stabilized at \qty{1}{\m} depth
and approximately \qty{1}{\m} from the lateral pool wall facing the fixed attachment point
(\cref{cable_modelization_fig_all_trajectories}). Surge and sway movements are conducted
individually to investigate the cable's response movement within or without its original vertical
plane. The robot repeats the same open-loop control in sway or surge with the same step profile one
after the other for each cable: 1) stabilisation with auto-depth at \qty{1}{\m} depth for
\qty{2}{\s}, 2) abrupt start with constant surge or sway command applied for \qty{2.5}{\s}, 3)
reversing of the thrusters for \qty{0.3}{\s} with a doubled velocity of opposite direction, to
obtain an abrupt stop, 4) slow down the thruster for half the initial velocity during \qty{0.2}{\s}
and 5) deactivation of the auto-depth control and complete stop of the thrusters after
\qty{13}{\s}. The robot then drifts slowly.

\begin{figure}
	\centering
	\inputtikz{figures/trajectories}
	\caption{
		Top view of typical trajectories of cable~\subref{cable_modelization_fig_image_weighted} end.
		Circles refer to initial positions and crosses refer to final positions
	}
	\label{cable_modelization_fig_all_trajectories}
\end{figure}

The starting position of the mini-ROV is manually determined with the help of a visual reference
point. It is not always perfectly aligned with the fixed point (see starting circles and
trajectories in \cref{cable_modelization_fig_all_trajectories}). To address this problem,
closed-loop control could be considered. However, conventional position sensors such as IMUs, DVLs
and USBLs are not suitable for an experiment in a compact metal tank, and the use of a visual
target limits the range of possible movements. Another option would be to use a traction bench.
However, such a setup would not be representative of the dynamics of the mini-ROV in interaction
with its tether. Nevertheless, the depth sensor allows for closed-loop depth control and the
desired direction of motion was dominant with its open-loop controller.

The experiment is repeated three times for each cable, under six experimental conditions, which
makes it possible to test the reactions of more or less taut cables with accelerations of different
amplitudes. Each experimental condition is a combination of a direction of motion (sway or surge),
a speed (\qty{0.3}{\m\per\s} or \qty{0.6}{\m\per\s}) and a starting point characterized by its
horizontal distance from the fixation point (about \qty{1.5}{\m} or \qty{2}{\m}).

\subsection{Results}

\subsubsection{General study}

\Cref{cable_modelization_fig_boxplot_all} represents the residuals
\cref{cable_modelization_eq_error_vt_definition} for the three model variants, namely standard
Catenary, \gammamodel{} and \thetagammamodel{}, for all sequences and all experimental conditions.
There is a general trend showing that the \thetagammamodel{} model has lower residuals than the
standard catenary model and the \gammamodel{} one. The \gammamodel{} variant is also more accurate
than the standard catenary model.

Regardless of the experimental condition, cable \subref{cable_modelization_fig_image_weighted} and
\subref{cable_modelization_fig_image_chain} show lower residuals for all three variants, as
expected regarding their physical properties (see
\cref{cable_modelization_tab_cables_characteristics}). Indeed, these cables are the closer to the
definition of the catenary because they are flexible and negatively buoyant. Sail rope
\subref{cable_modelization_fig_image_rope} and elastic rope
\subref{cable_modelization_fig_image_elastic} are also negatively buoyant but lighter with regards
to their volume so their residual is higher. Indeed, the drag forces damping effects wins over the
weight. The floating cable~\subref{cable_modelization_fig_image_floating} behaves like an inverted
catenary. It has a high residual for the standard catenary model (2\textsuperscript{nd} worst
median) but shows a good fit for the \thetagammamodel{} model (3\textsuperscript{rd} best median).

For the stiffer
cables~\subref{cable_modelization_fig_image_coaxial},~\subref{cable_modelization_fig_image_four_pairs}
and~\subref{cable_modelization_fig_image_two_pairs}, the improvement in accuracy for the
\thetagammamodel{} variant still exists, albeit to a lesser extent. Indeed, the stiffer the cable,
the higher the residuals. These cables have permanent deformations as a result of being coiled.
Furthermore they are neutrally buoyant. As a result, the impact of weight is insufficient to
generate the catenary shape.

The section below examines the three model variants in greater details, distinguishing between the
experiments involving predominant surge and sway motion.
Cables~\subref{cable_modelization_fig_image_two_pairs} is chosen as representative of the stiff
cables with neutral buoyancy and cable~\subref{cable_modelization_fig_image_weighted} is chosen
among the flexible cables that are negatively buoyant (as defined
in~\cref{cable_modelization_subsec_cable}).

\begin{figure}
	\centering
	\inputtikz{figures/cables_boxplot}
	\caption{
		Box plots of  all 8 cables under all experimental parameters (speed, direction, distance) with the median,
		1\textsuperscript{st} and 3\textsuperscript{rd} quartile, the 1\textsuperscript{st} and 99\textsuperscript{th} percentile
		of model residuals $\pdistance{v}$ (blue), $\pdistance{\gamma}$ (red) and $\pdistance{\theta\gamma}$ (gray).
	}
	\label{cable_modelization_fig_boxplot_all}
\end{figure}

\subsubsection{Detailed analysis of cables~\subref{cable_modelization_fig_image_two_pairs} and~\subref{cable_modelization_fig_image_weighted} as group representatives}

\begin{figure}
	\centering
	\inputtikz{figures/cable3_boxplot}
	\caption{
		Box plots of the two-pair network cable, cable~\subref{cable_modelization_fig_image_two_pairs} residuals
		$\pdistance{v}$ in blue,
		$\pdistance{\gamma}$ in red and $\pdistance{\theta\gamma}$ in gray
	}
	\label{cable_modelization_tab_parameters_cable3}
\end{figure}

\begin{figure}
	\centering
	\inputtikz{figures/cable6_boxplot}
	\caption{
		Box plots of the leaded seamstress rope, cable~\subref{cable_modelization_fig_image_weighted} residuals
		$\pdistance{v}$ in blue,
		$\pdistance{\gamma}$ in red and $\pdistance{\theta\gamma}$ in gray
	}
	\label{cable_modelization_tab_parameters_cable6}
\end{figure}

\cref{cable_modelization_tab_parameters_cable3} and~\cref{cable_modelization_tab_parameters_cable6}
show the box plots of the model residuals for cable~\subref{cable_modelization_fig_image_two_pairs}
and~\subref{cable_modelization_fig_image_weighted} according to experimental parameters. First, the
sway movements are displayed for a distance of \qty{2}{\m} and two different velocities. Then the
sway movements are detailed for two different velocities and two starting points. These Tables show
that the improvements of accuracy of the \thetagammamodel{} model is consistent across experimental
parameters.

For both groups representatives, the residuals of the models are smaller for slower movements and
greater distance from the attachment point (\textit{i.e.} more tension in the cable). The
\gammamodel{} model displays a higher accuracy for lateral sway movements compared to forward surge
movements. Conversely, our \thetagammamodel{} is significantly better for surge movements in
comparison of the \gammamodel{}. The gain in accuracy is even greater for
cable~\subref{cable_modelization_fig_image_weighted}'s group.

\subsubsection{Detailed analysis for experimental conditions of 0.6~m~s\textsuperscript{-1} speed and 2~m distance for cables~\subref{cable_modelization_fig_image_two_pairs}
	and~\subref{cable_modelization_fig_image_weighted} as group representatives}

\begin{figure}
	\centering
	\inputtikz{figures/sway_3}
	\caption{
		Model residuals, robot speed and $\theta$, $\gamma$ angles over time. Speed command of \qty{0.6}{\m\per\s} and
		\qty{2}{\m} distance with cable~\subref{cable_modelization_fig_image_two_pairs}, sway motion
	}
	\label{cable_modelization_fig_sequences_cable3sway}
\end{figure}

\begin{figure}
	\centering
	\inputtikz{figures/sway_6}
	\caption{
		Model residuals, robot speed and $\theta$, $\gamma$ angles over time. Speed command of \qty{0.6}{\m\per\s} and
		\qty{2}{\m} distance with cable~\subref{cable_modelization_fig_image_weighted}, sway motion
	}
	\label{cable_modelization_fig_sequences_cable6sway}
\end{figure}

\begin{figure}
	\centering
	\inputtikz{figures/surge_3}
	\caption{
		Model residuals, robot speed and $\theta$, $\gamma$ angles over time. Speed command of \qty{0.6}{\m\per\s} and
		\qty{2}{\m} distance with cable~\subref{cable_modelization_fig_image_two_pairs}, surge motion
	}
	\label{cable_modelization_fig_sequences_cable3surge}
\end{figure}

\begin{figure}
	\centering
	\inputtikz{figures/surge_6}
	\caption{
		Model residuals, robot speed and $\theta$, $\gamma$ angles over time. Speed command of \qty{0.6}{\m\per\s} and
		\qty{2}{\m} distance. with cable~\subref{cable_modelization_fig_image_weighted}, surge motion
	}
	\label{cable_modelization_fig_sequences_cable6surge}
\end{figure}

\Cref{cable_modelization_fig_sequences_cable3surge,cable_modelization_fig_sequences_cable3sway,cable_modelization_fig_sequences_cable6surge,cable_modelization_fig_sequences_cable6sway}
shows the three model variants residuals, angles $\theta$ and $\gamma$, as well as the robot's
speed. These figures show in detail the different phases of the movement: an initial stable
position (hydrostatic equilibrium phase), an abrupt start (dynamic phase), a phase of continuous
application of a constant command (steady state phase), an abrupt stop (dynamic phase) and a
subsequent hydrostatic phase.

A clear correlation appears between the robot's dynamics and the angle $\gamma$ for sway motion.
The estimated $\gamma$ angle consistently increases with speed during sway movements, whereas
$\theta$ remains close to zero (\cref{cable_modelization_fig_sequences_cable3sway}
and~\cref{cable_modelization_fig_sequences_cable6sway}). As anticipated, surge movements have a
noticeable impact on cable~\subref{cable_modelization_fig_image_two_pairs}'s estimated angle
$\theta$, with $\gamma$ maintaining a small value
(\cref{cable_modelization_fig_sequences_cable3surge}). For
cable~\subref{cable_modelization_fig_image_weighted}, both angles increase with velocity, with
$\theta$ being greater than the one observed for sway movement
(\cref{cable_modelization_fig_sequences_cable6surge}). However, since $\gamma$ is larger than
expected, let us examine the trajectory displayed in \cref{cable_modelization_fig_all_trajectories}
relative to surge motion with \qty{0.6}{\m\per\s} and \qty{2}{\m}. While the primary motion is
surge, the trajectory of the robot is not precisely aligned with the fixed point causing the cable
to move laterally. Although this effect is present throughout all sequences, it is particularly
pronounced in this trajectory which explains the values for $\gamma$. This demonstrates that our
\thetagammamodel{} model is capable of simultaneously managing significant sway and surge
movements. Finally, one can clearly see that cable~\subref{cable_modelization_fig_image_weighted}
restores balance towards the vertical plane and the standard model at a faster rate compared to the
lighter cable~\subref{cable_modelization_fig_image_two_pairs}. As a result, physics-based modelling
of angle dynamics opens up promising perspectives for model-based control synthesis.

\section{Conclusion}
\label{cable_modelization_sec_conclusion}

This work introduced an augmented catenary model that accounts for the hydrodynamic effects on the
tether when the ROV performs surge and sway motion. It incorporates the models presented in the
state of the art, which can be found by zeroing one or both of the angles proposed as augmentation
parameters.

The results drawn from the experiments carried out on eight different cables show that the
augmented model provides a better estimate of the shape of all cables during dynamic phases,
compared with the standard catenary model. The accuracy is best for flexible cables with adequate
negative buoyancy. In addition, cables with the lowest residuals relative to the standard catenary
model lead to the best improvements in accuracy using the augmented model. Their measured shape
features greater planarity and homogenous tilting of their support plane during motion.

In light of these results, the augmented catenary model will be used in model-based controllers to
estimate cable shape and its lowest position for deployment scenarios of tethered robots.
